\clearpage
\fancyhead[L]{\small COMET.}
\newcommand{\CometFourSevenTwoDecl}{\hbox{\vrule width0pt height8bp depth0pt\special{pdf:bcontent}\special{pdf:code q 0 g 4.909091 8.000000 m 4.000000 7.090909 l 4.000000 6.909091 l 3.818182 6.727273 l 3.818182 6.545455 l 1.818182 4.545455 l 1.454545 4.545455 l 1.272727 4.363636 l 1.090909 4.363636 l 0.363636 3.636364 l 0.363636 3.454545 l 0.181818 3.272727 l 0.181818 3.090909 l 0.000000 2.909091 l 0.000000 1.636364 l 0.181818 1.454545 l 0.181818 1.272727 l 1.090909 0.363636 l 1.818182 0.363636 l 2.000000 0.181818 l 2.181818 0.181818 l 2.363636 0.363636 l 3.272727 0.363636 l 3.636364 0.727273 l 3.818182 0.727273 l 4.181818 1.090909 l 4.181818 1.818182 l 4.363636 2.000000 l 4.363636 2.181818 l 4.181818 2.363636 l 4.181818 2.545455 l 4.000000 2.727273 l 4.000000 3.090909 l 4.181818 3.272727 l 4.363636 3.272727 l 4.909091 3.818182 l 5.272727 3.818182 l 5.454545 4.000000 l 5.636364 4.000000 l 6.000000 4.363636 l 6.000000 4.545455 l 6.363636 4.909091 l 6.545455 4.909091 l 6.727273 5.090909 l 6.909091 5.090909 l 7.818182 6.000000 l 8.000000 6.000000 l 8.000000 6.181818 l 8.181818 6.363636 l 8.363636 6.363636 l 8.545455 6.545455 l 8.545455 6.727273 l 8.727273 6.909091 l 8.545455 7.090909 l 8.727273 7.272727 l 8.545455 7.454545 l 8.000000 7.454545 l 7.636364 7.818182 l 6.545455 7.818182 l 6.363636 7.636364 l 6.181818 7.818182 l 6.000000 7.818182 l 5.818182 8.000000 l h 1.454545 3.636364 m 1.636364 3.818182 l 2.363636 3.818182 l 2.545455 3.636364 l 2.909091 3.636364 l 3.272727 3.272727 l 3.272727 2.909091 l 3.454545 2.727273 l 3.454545 2.181818 l 3.272727 2.000000 l 3.272727 1.818182 l 3.090909 1.636364 l 3.090909 1.454545 l 2.545455 0.909091 l 1.636364 0.909091 l 1.454545 1.090909 l 1.272727 1.090909 l 0.727273 1.636364 l 0.727273 3.090909 l 1.272727 3.636364 l h 3.272727 5.090909 m 3.454545 5.272727 l 3.636364 5.272727 l 3.818182 5.090909 l 3.818182 4.909091 l 3.272727 4.363636 l 2.909091 4.363636 l 2.727273 4.545455 l 2.727273 4.727273 l 3.090909 5.090909 l h 7.090909 7.272727 m 7.454545 7.272727 l 7.636364 7.090909 l 4.727273 4.181818 l 4.545455 4.181818 l 3.818182 3.454545 l 3.636364 3.454545 l 3.272727 3.818182 l 3.818182 4.363636 l 4.000000 4.363636 l 4.545455 4.909091 l 4.727273 4.909091 l 4.909091 5.090909 l 4.909091 5.272727 l 5.272727 5.636364 l 5.454545 5.636364 l h 5.454545 7.272727 m 5.818182 7.272727 l 6.000000 7.090909 l 4.727273 5.818182 l 4.545455 5.818182 l 4.363636 5.636364 l 4.181818 5.636364 l 4.000000 5.818182 l h f* Q}\special{pdf:econtent}\kern8.909091bp}}
\newcommand{\CometFourSevenTwoAR}{\hbox{\vrule width0pt height8bp depth0pt\special{pdf:bcontent}\special{pdf:code q 0 g 5.454545 8.000000 m 5.272727 7.818182 l 4.909091 7.818182 l 4.363636 7.272727 l 4.363636 7.090909 l 1.636364 4.363636 l 1.272727 4.363636 l 1.090909 4.181818 l 0.909091 4.181818 l 0.181818 3.454545 l 0.181818 2.909091 l 0.000000 2.727273 l 0.000000 1.636364 l 0.181818 1.454545 l 0.181818 1.272727 l 1.090909 0.363636 l 1.454545 0.363636 l 1.636364 0.181818 l 2.545455 0.181818 l 2.727273 0.363636 l 3.090909 0.363636 l 3.272727 0.545455 l 3.454545 0.545455 l 4.000000 1.090909 l 4.000000 1.636364 l 4.181818 1.818182 l 4.181818 2.363636 l 4.000000 2.545455 l 4.000000 2.909091 l 4.545455 3.454545 l 4.727273 3.454545 l 5.090909 3.818182 l 5.454545 3.818182 l 6.000000 4.363636 l 6.181818 4.363636 l 6.363636 4.545455 l 6.363636 4.727273 l 6.909091 5.272727 l 7.090909 5.272727 l 7.272727 5.454545 l 7.454545 5.454545 l 7.454545 5.636364 l 8.363636 6.545455 l 8.545455 6.545455 l 8.727273 6.727273 l 8.727273 7.272727 l 8.363636 7.636364 l 8.181818 7.636364 l 8.000000 7.454545 l 7.818182 7.454545 l 7.454545 7.818182 l 6.727273 7.818182 l 4.909091 6.000000 l 4.909091 5.818182 l 4.363636 5.272727 l 4.181818 5.272727 l 3.454545 4.545455 l 3.272727 4.545455 l 2.909091 4.181818 l 2.727273 4.181818 l 2.545455 4.363636 l 2.545455 4.545455 l 3.272727 5.272727 l 3.454545 5.272727 l 4.545455 6.363636 l 4.545455 6.545455 l 5.272727 7.272727 l 5.454545 7.272727 l 5.636364 7.454545 l 5.818182 7.454545 l 6.000000 7.636364 l 6.000000 7.818182 l 5.818182 8.000000 l h 1.636364 3.636364 m 1.818182 3.818182 l 2.545455 3.818182 l 3.454545 2.909091 l 3.454545 2.181818 l 3.272727 2.000000 l 3.272727 1.636364 l 3.090909 1.454545 l 3.090909 1.272727 l 2.727273 0.909091 l 1.454545 0.909091 l 0.545455 1.818182 l 0.545455 2.545455 l 0.727273 2.727273 l 0.727273 2.909091 l 1.454545 3.636364 l h 7.090909 7.090909 m 7.272727 7.272727 l 7.454545 7.272727 l 7.636364 7.090909 l 7.636364 6.909091 l 6.000000 5.272727 l 5.818182 5.272727 l 4.727273 4.181818 l 4.545455 4.181818 l 4.181818 3.818182 l 4.181818 3.636364 l 3.818182 3.272727 l 3.636364 3.272727 l 3.272727 3.636364 l 5.272727 5.636364 l 5.454545 5.636364 l 6.909091 7.090909 l h f* Q}\special{pdf:econtent}\kern8.909091bp}}

\hypertarget{native-471-lower}{}
\begin{center}
\rule{0.55\linewidth}{0.4pt}\\[2pt]
Astronomische Nachrichten. Band XXII. Nr. 511. Seite 111. 112. 1844 August 1.\\[2pt]
\rule{0.55\linewidth}{0.4pt}
\end{center}
\begin{center}Beobachtungen und Elemente des \emph{Cometen.}\end{center}

Herr Hofrath \textsc{Gauss} hat auf der Göttinger Sternwarte bei ganz günstiger Luft, folgende Beobachtung des Cometen erhalten:

\Needspace{5\baselineskip}
\SourceLeaf{472}
\begin{center}
\setlength{\tabcolsep}{3pt}
\begin{tabular}{llll}
1844 Juli 23. & $10^{\rm u}\ 46^{\rm m}\ 32^{\rm s}$ m. Z. & Decl. \CometFourSevenTwoDecl & $=+36^\circ\ 29'\ 51''2$ \\
 & $10\ 56\ 4$\quad --- & AR. \CometFourSevenTwoAR & $=226\ 7\ 13.8$ \\
\end{tabular}
\end{center}

Die Gerade Aufsteigung beruhet auf 4 Vergleichungen, die Declination nur auf Einer.

\bigskip
\hypertarget{native-472-notice}{}
\begin{center}
\rule{0.55\linewidth}{0.4pt}\\[2pt]
Astronomische Nachrichten. Band XXII. Nr. 512. Seite 115. 116. 1844 August 17.\\[2pt]
\rule{0.55\linewidth}{0.4pt}
\end{center}

[Mittheilung der von Dr. \textsc{Goldschmidt} berechneten Elemente des Cometen.]

\clearpage
\hypertarget{native-472-august}{}
\begin{center}
\rule{0.55\linewidth}{0.4pt}\\[2pt]
Astronomische Nachrichten. Band XXII. Nr. 513. Seite 141..144. 1844 August 22.\\[2pt]
\rule{0.55\linewidth}{0.4pt}
\end{center}
\begin{flushright}Göttingen, 1844 Aug. 8.\end{flushright}

Seit dem 23. v. M. habe ich zwar an mehreren Abenden den Cometen auf kurze Zeit gesehen, aber erst gestern wieder eine Beobachtung erhalten können; aber auch schon bei der dritten Vergleichung kamen wieder Wolken, die weitere Beobachtungen unmöglich machten. Was ich erhalten habe, ist folgendes, wobei zwei Sterne aus \textsc{Bessels} Zone Nr. 364 zur Vergleichung benutzt wurden.

\begin{center}
\setlength{\tabcolsep}{3pt}
\begin{tabular}{llll}
1844 Aug. 7. & $9^{\rm u}\ 16^{\rm m}\ 29^{\rm s}$ M. Z. & A. R. & $=213^\circ\ 52'\ 6''8$ \\
 & $9\ 24\ 48\mathord{\cdot}$\quad --- & Decl. & $=+24\ 50\ 21.5$ \\
\end{tabular}
\end{center}

{\rightskip=0pt plus 1em Nach einer flüchtigen Ueberschlagung geben Dr. \textsc{Goldschmidts} parabolische Elemente die G. A. um $10''$ grösser, die Abweichung um $16''$ kleiner.\par}

\begin{flushright}Göttingen, 1844. Aug. 10.\end{flushright}

Ich übersende Ihnen abermals eine Ortsbestimmung des Cometen, die ich gestern Abend gemacht habe. Zur Vergleichung dienten 2 Sterne aus \textsc{Bessels} Zone 364, der erste auch aus Zone 460.

\begin{center}
\setlength{\tabcolsep}{3pt}
\begin{tabular}{llll}
1844 Aug. 9. & $9^{\rm u}\ 36^{\rm m}\ 54^{\rm s}$ M. Z. & G. R. & $212^\circ\ 40'\ 23''8$ \\
 & $9\ 35\ 8$\quad --- & Abw. & $23\ 18\ 16.6$ \\
\end{tabular}
\end{center}

\textsc{Goldschmidts} Elemente geben die G. A. $5''$ grösser, die Declination $16''$ kleiner.
